% GENERATED FILE. Edit canonical lessons, then run npm run generate:books.
% canonical-source-hash: fnv1a64:22ac91b8efd6e1e1
% canonical-lessons: BN-C482-jetha, BN-C482-meso, BN-C482-bhagne, BN-C482-bhaijhi, BN-C482-beyai

\chapter{More Relatives}
\label{ch:bn-more-relatives}
\hlchaptermodality{482}

\begin{chapteropening}
\textbf{By the end of this chapter:} \emph{I can name more uncles, nephews, nieces and new relations.}

Complete the last lesson of chapter 482: I can name more uncles, nephews, nieces and new relations.
\end{chapteropening}

\section[\texorpdfstring{\bn{জেঠা} (jeṭhā)}{জেঠা (jeṭhā)}]{\bn{জেঠা} (jeṭhā) — a father's elder brother}

\label{lesson:BN-C482-jetha}



\begin{quote}
Before the new one: say the Bengali for an elder brother's wife, then the Bengali for a paternal grandmother.
\end{quote}

\subsection*{You'll want to know: \bn{জেঠা}}
\textbf{\bn{জেঠা}} — \emph{jeṭhā} — \textquotedblleft{}a father's elder brother\textquotedblright{}.

\begin{quote}
\bn{জেঠা} \bn{বাবার} \bn{চেয়ে} \bn{পাঁচ} \bn{বছরের} \bn{বড়।}
\end{quote}

\begin{quote}
\emph{jeṭhā bābār cheye pā̃ch bôchhôrer bôṛo.}
\end{quote}

\begin{quote}
My father's elder brother is five years older than my father.
\end{quote}

\bn{কাকা} (taught) is the father's younger brother; \bn{জেঠা} the elder.

\begin{cousinweb}
\textbf{\bn{জেঠা}} (\emph{jeṭhā}) is listed by CDIAL for Bengali under Sanskrit \emph{jyēṣṭhatāta}, \textquotedblleft{}eldest father\textquotedblright{} (CDIAL 5289).
\end{cousinweb}

\begin{grammarlens}[title={What you've built}]
One more word for describing what you lived and felt.
\end{grammarlens}

\subsection*{Your turn}
\begin{itemize}
\raggedright
  \item \emph{Say it:} \emph{jeṭhā}
  \item \emph{Say it:} \emph{jeṭhā}, once more
  \item \emph{Say it:} \emph{jeṭhā bābār cheye pā̃ch bôchhôrer bôṛo.} — My father's elder brother is five years older than my father.
  \item \emph{Recall:} say the Bengali for an elder brother's wife, then the Bengali for a paternal grandmother, then say \emph{jeṭhā} again
\end{itemize}

\begin{tcolorbox}[breakable,colback=teal!4,colframe=teal!35!black,title={Before you move on}]
What does \bn{জেঠা} mean? (\textquotedblleft{}A father's elder brother\textquotedblright{}.) Say it once more.
\end{tcolorbox}
\section[\texorpdfstring{\bn{মেসো} (meso)}{মেসো (meso)}]{\bn{মেসো} (meso) — a maternal aunt's husband}

\label{lesson:BN-C482-meso}



\begin{quote}
Before the new one: say the Bengali for a paternal grandmother, then the Bengali for a father's elder brother.
\end{quote}

\subsection*{You'll want to know: \bn{মেসো}}
\textbf{\bn{মেসো}} — \emph{meso} — \textquotedblleft{}a maternal aunt's husband\textquotedblright{}.

\begin{quote}
\bn{মেসো} \bn{আমাকে} \bn{একটা} \bn{নতুন} \bn{সাইকেল} \bn{কিনে} \bn{দিলেন।}
\end{quote}

\begin{quote}
\emph{meso āmāke ekṭā nôtun sāikel kine dilen.}
\end{quote}

\begin{quote}
My aunt's husband bought me a new bicycle.
\end{quote}

\begin{cousinweb}
\textbf{\bn{মেসো}} (\emph{meso}) is listed by CDIAL for Bengali, as a mother's sister's husband, under Sanskrit \emph{mātuḥṣvasṛ}, \textquotedblleft{}mother's sister\textquotedblright{} (CDIAL 10001); \bn{মাসি} (taught) is the aunt herself.
\end{cousinweb}

\begin{grammarlens}[title={What you've built}]
One more word for describing what you lived and felt.
\end{grammarlens}

\subsection*{Your turn}
\begin{itemize}
\raggedright
  \item \emph{Say it:} \emph{meso}
  \item \emph{Say it:} \emph{meso}, once more
  \item \emph{Say it:} \emph{meso āmāke ekṭā nôtun sāikel kine dilen.} — My aunt's husband bought me a new bicycle.
  \item \emph{Recall:} say the Bengali for a paternal grandmother, then the Bengali for a father's elder brother, then say \emph{meso} again
\end{itemize}

\begin{tcolorbox}[breakable,colback=teal!4,colframe=teal!35!black,title={Before you move on}]
What does \bn{মেসো} mean? (\textquotedblleft{}A maternal aunt's husband\textquotedblright{}.) Say it once more.
\end{tcolorbox}
\section[\texorpdfstring{\bn{ভাগ্নে} (bhāgne)}{ভাগ্নে (bhāgne)}]{\bn{ভাগ্নে} (bhāgne) — a sister's son; a nephew}

\label{lesson:BN-C482-bhagne}



\begin{quote}
Before the new one: say the Bengali for a father's elder brother, then the Bengali for a maternal aunt's husband.
\end{quote}

\subsection*{You'll want to know: \bn{ভাগ্নে}}
\textbf{\bn{ভাগ্নে}} — \emph{bhāgne} — \textquotedblleft{}a sister's son; a nephew\textquotedblright{}.

\begin{quote}
\bn{আমার} \bn{ছোট} \bn{ভাগ্নের} \bn{ক্রিকেট} \bn{খেলতে} \bn{খুব} \bn{ভালো} \bn{লাগে।}
\end{quote}

\begin{quote}
\emph{āmār chhoṭo bhāgner krikeṭ khelte khub bhālo lāge.}
\end{quote}

\begin{quote}
My little nephew really likes playing cricket.
\end{quote}

\begin{cousinweb}
\textbf{\bn{ভাগ্নে}} (\emph{bhāgne}) comes from Sanskrit \emph{bhāgineya}, \textquotedblleft{}a sister's son\textquotedblright{} (CDIAL 9433), worn down in Bengali.
\end{cousinweb}

\begin{grammarlens}[title={What you've built}]
One more word for describing what you lived and felt.
\end{grammarlens}

\subsection*{Your turn}
\begin{itemize}
\raggedright
  \item \emph{Say it:} \emph{bhāgne}
  \item \emph{Say it:} \emph{bhāgne}, once more
  \item \emph{Say it:} \emph{āmār chhoṭo bhāgner krikeṭ khelte khub bhālo lāge.} — My little nephew really likes playing cricket.
  \item \emph{Recall:} say the Bengali for a father's elder brother, then the Bengali for a maternal aunt's husband, then say \emph{bhāgne} again
\end{itemize}

\begin{tcolorbox}[breakable,colback=teal!4,colframe=teal!35!black,title={Before you move on}]
What does \bn{ভাগ্নে} mean? (\textquotedblleft{}A sister's son; a nephew\textquotedblright{}.) Say it once more.
\end{tcolorbox}
\section[\texorpdfstring{\bn{ভাইঝি} (bhāijhi)}{ভাইঝি (bhāijhi)}]{\bn{ভাইঝি} (bhāijhi) — a brother's daughter; a niece}

\label{lesson:BN-C482-bhaijhi}



\begin{quote}
Before the new one: say the Bengali for a maternal aunt's husband, then the Bengali for a sister's son.
\end{quote}

\subsection*{You'll want to know: \bn{ভাইঝি}}
\textbf{\bn{ভাইঝি}} — \emph{bhāijhi} — \textquotedblleft{}a brother's daughter; a niece\textquotedblright{}.

\begin{quote}
\bn{ভাইঝির} \bn{জন্মদিনে} \bn{আমি} \bn{একটা} \bn{পুতুল} \bn{কিনলাম।}
\end{quote}

\begin{quote}
\emph{bhāijhir jônmôdine āmi ekṭā putul kinlām.}
\end{quote}

\begin{quote}
I bought a doll for my niece's birthday.
\end{quote}

\begin{cousinweb}
\textbf{\bn{ভাইঝি}} (\emph{bhāijhi}) puts \bn{ভাই} (taught), brother, with \emph{jhi}, \textquotedblleft{}daughter\textquotedblright{}, which CDIAL lists for Bengali under Sanskrit \emph{duhitṛ}, \textquotedblleft{}daughter\textquotedblright{} (CDIAL 6481).
\end{cousinweb}

\begin{grammarlens}[title={What you've built}]
One more word for describing what you lived and felt.
\end{grammarlens}

\subsection*{Your turn}
\begin{itemize}
\raggedright
  \item \emph{Say it:} \emph{bhāijhi}
  \item \emph{Say it:} \emph{bhāijhi}, once more
  \item \emph{Say it:} \emph{bhāijhir jônmôdine āmi ekṭā putul kinlām.} — I bought a doll for my niece's birthday.
  \item \emph{Recall:} say the Bengali for a maternal aunt's husband, then the Bengali for a sister's son, then say \emph{bhāijhi} again
\end{itemize}

\begin{tcolorbox}[breakable,colback=teal!4,colframe=teal!35!black,title={Before you move on}]
What does \bn{ভাইঝি} mean? (\textquotedblleft{}A brother's daughter; a niece\textquotedblright{}.) Say it once more.
\end{tcolorbox}
\section[\texorpdfstring{\bn{বেয়াই} (beyāi)}{বেয়াই (beyāi)}]{\bn{বেয়াই} (beyāi) — a child's father-in-law; a co-father-in-law}

\label{lesson:BN-C482-beyai}



\begin{quote}
Before the new one: say the Bengali for a sister's son, then the Bengali for a brother's daughter.
\end{quote}

\subsection*{You'll want to know: \bn{বেয়াই}}
\textbf{\bn{বেয়াই}} — \emph{beyāi} — \textquotedblleft{}a child's father-in-law; a co-father-in-law\textquotedblright{}.

\begin{quote}
\bn{বিয়ের} \bn{পরে} \bn{দুই} \bn{বেয়াই} \bn{বন্ধু} \bn{হলেন।}
\end{quote}

\begin{quote}
\emph{biyer pôre dui beyāi bôndhu hôlen.}
\end{quote}

\begin{quote}
After the wedding the two fathers-in-law became friends.
\end{quote}

\begin{cousinweb}
\textbf{\bn{বেয়াই}} (\emph{beyāi}) is listed by CDIAL for Bengali, as a son's or daughter's father-in-law, under Sanskrit \emph{vivāhin} (CDIAL 11925), from the same root as \bn{বিয়ে} (taught), a wedding.
\end{cousinweb}

\begin{grammarlens}[title={What you've built}]
One more word for describing what you lived and felt.
\end{grammarlens}

\subsection*{Your turn}
\begin{itemize}
\raggedright
  \item \emph{Say it:} \emph{beyāi}
  \item \emph{Say it:} \emph{beyāi}, once more
  \item \emph{Say it:} \emph{biyer pôre dui beyāi bôndhu hôlen.} — After the wedding the two fathers-in-law became friends.
  \item \emph{Recall:} say the Bengali for a sister's son, then the Bengali for a brother's daughter, then say \emph{beyāi} again
\end{itemize}

\begin{tcolorbox}[breakable,colback=teal!4,colframe=teal!35!black,title={Before you move on}]
What does \bn{বেয়াই} mean? (\textquotedblleft{}A child's father-in-law; a co-father-in-law\textquotedblright{}.) Say it once more.
\end{tcolorbox}
